\documentclass{article}

\usepackage{notes}

\title{GRT Seminar Spring 2026 -- Potent Categorical Representations}
\author{Notes by John S.\ Nolan, speakers listed below}

\begin{document}

\maketitle

\begin{abstract}
	This semester, the GRT Seminar will focus on potent categorical representations, following the 2026 ICM address of Nadler--Ben-Zvi (based on joint work with Stefanich).
  There are many open questions!
\end{abstract}

\tableofcontents

\section{(1/22) David Nadler -- Introduction}

We're going to start with the first page or so of the paper.
The rough idea here is to instantiate some of the ideas in Constantin Teleman's ICM address.

\subsection{Categorical representations}

Our motivation for studying categorical representations comes from the following analogy.

If $\Gamma$ is a finite group, then we look at the group ring $\CC[\Gamma]$ 
We may think of this as a ring of distributions, with multiplication given by convolution.
Studying representations of $\Gamma$ corresponds to studying modules over $\CC[\Gamma]$.

In general, if $\Cc$ is a symmetric monoidal ($\infty$-)category, e.g.\ $(\Vect_\CC, \otimes)$, we may consider an algebra object $A \in \Cc$.
We can study the category of modules $\Mod_A(\Cc)$.

We can specialize this very general framework to give a notion of categorical representations.
Let $G$ be a connected reductive group over $\CC$.

\begin{dfn}
  A \emph{categorical $G$-representation} is a (presentable) stable $\CC$-linear $\infty$-category which is a module over $D(G)$, the derived $\infty$-category of $D$-modules on $G$.
  We'll write $G\dCat$ for the $\infty$-category of categorical $G$-representations.
\end{dfn}

\begin{rmk}
  Morally, the reason we use $D$-modules here is because they quantize the cotangent bundle.
  The really important object is $T^*G$ (with two moment maps $T^*G \rightrightarrows\gfr^*$ corresponding to the left and right actions of $G$).
  Everything here is secretly symplectic!
\end{rmk}

\begin{ex}
  Let $G \curvearrowright X$.
  Then we get an action $D(G) \curvearrowright D(X)$.
  This really corresponds to the moment map $T^* X \to \gfr^*$.
\end{ex}

\subsection{Categorical representations of $\GG_m$}

Let's see what happens when $G = \GG_m$.

\begin{thm}
  There is a monoidal equivalence $D(\GG_m) \simeq \QC(\AA^1 / \ZZ)$.
\end{thm}

Indeed, recall that a $D$-module on $\GG_m$ comes with two pieces of data:
\begin{itemize}
  \item An action of the Euler vector field $t \partial_t$, corresponding to the coordinate on $\AA^1$.
  \item An action of $t$ and $t\inv$, corresponding to translation-equivariance.
\end{itemize}
This correspondence sends the $\delta$ distribution at $1$ to the structure sheaf.

\begin{cor}
  There is an equivalence $\GG_m\dCat \simeq \QC(\AA^1 / \ZZ)\dMod$.
\end{cor}

Reasonable objects within $\QC(\AA^1 / \ZZ)\dMod$ have ``singular support'' in (the 2-shifted version of) $T^*(\AA^1 / \ZZ)$.
This translates into a notion of ``position'' and ``momentum'' for such objects.

\begin{ex}
  If we consider a point $\lambda \in \AA^1 / \ZZ$, the module category $\QC(\lambda)$ corresponds to
  \[
    D_\lambda(\GG_m) = \angles{\Osc_{\GG_m}, d - \frac{\lambda}{z} dz} \subset \Dsf(\GG_m).
  \]
  Roughly: the position of a $\QC(\AA^1 / \ZZ)$-module corresponds to the monodromy.
  This relates to Langlands parameters.
\end{ex}

Momentum should correspond to ``equivariant parameters.''
To capture this, we must work microlocally.
This is the motivation for considering potent categorical representations.

\subsection{Another motivation}

Sufficiently dualizable $2$-categories give rise to 3d TQFTs.
One can obtain some 3d TQFTs by taking $S^1$-compactifications of 4d TQFTs, i.e.\ setting $Z_3(\pt) = Z_4(S^1)$.
If $Z_4$ is pure gauge theory, $Z_3$ corresponds to considering gauge theory with classical fields $G / G$.
Remember that $G / G$ can be thought of as the free loop space $\Lsc BG$.
The ``equivariant parameters'' here live in the numerator $G$, while the ``monodromies'' / ``$G$-representations'' come from the denominator $G$.

\subsection{Defining potent categorical representations}

We can write $G\dCat = 2\IndCoh(G_\dR)$.
Thus a possible extension is:

\begin{dfn}
  Let $G\dCat^\pot = 2\IndCoh(G / G)^\Tate$.
\end{dfn}

This contains $G\dCat$ as the full $2$-subcategory set-theoretically supported at $1 \in G$.
There are other subtleties (e.g.\ shearing, genuine equivariance) that should appear. 
We end up with a theorem:

\begin{thm}
  There is an equivalence $\GG_m\dCat^\pot \simeq 2\IndCoh(\AA^1 / \ZZ)$.
\end{thm}

We may also define a Lie algebra version:

\begin{dfn}
  Let $G\dCat^{\pot,\gr} = 2\IndCoh(\gfr / G)^\Tate$.
\end{dfn}

Our goal this semester is to understand these definitions and theorems.

\section{(1/29) Samuel Mayonnaise -- Loop Spaces and $D$-Modules}

Our goal for today is to prove the following.

\begin{thm}
  Let $X$ be a smooth scheme.
  Then there is an equivalences
  \[
    D(X) \otimes k[u, u\inv] \simeq \IndCoh(\Lc X)^\Tate.
  \]
\end{thm}

The ``Tate construction'' here captures an $S^1$-localization.
To understand the above statement, we will need to recall some relevant definitions.

\subsection{Loop stacks}

Let $X$ be a derived stack, i.e.\ a (covariant) functor from derived commutative rings to spaces satisfying a sheaf condition.

\begin{rmk}
  There are many models for ``derived commutative rings'' that we can use.
  In characteristic zero we may model derived commutative rings by connective commutative differential graded algebras.
  In general we can work with simplicial commutative rings.
\end{rmk}

The \emph{loop stack} of $X$ is $\Lc X = \ul{\Hom}(S^1, X)$.
Here $S^1$ is viewed as a derived stack by sheafifying the constant sheaf with value $S^1$.
More generally, sheafifying constant sheaves gives a definition of the \emph{Betti stack} of any topological space.

\begin{rmk}
  The simplicial perspective can be useful in understanding this.
  We may think of the Betti stack of $S^1$ as a coequalizer of $\pt \rightrightarrows \pt$.
  This corresponds to building up $S^1$ as a simplicial set with nondegenerate $0$- and $1$-simplices.
\end{rmk}

\begin{ex}
  If $X$ is a Betti stack, then $\Lc X$ is the Betti stack corresponding to the (topological) loop space of $X$.
\end{ex}

Let's think of $S^1$ as a (homotopy) pushout:
\[
  \begin{tikzcd}
    S^0 \rar \dar & \pt \dar \\
    \pt \rar & S^1.
  \end{tikzcd}
\]
This allows us to write
\[
  \Lc X = \ul{\Hom}(\pt \sqcup_{S^1} \pt, X) = \ul{\Hom}(\pt, X) \times_{\ul{\Hom}(S^0, X)} \ul{\Hom}(\pt, X) = X \times_{X \times X} X. 
\]
Thus we may interpret $\Lc X$ as the derived self-intersection of the diagonal of $X$.
This suggests that $D$-modules may potentially appear, since the normal bundle of the diagonal is the tangent bundle of $X$.

Assuming $X$ has affine diagonal, we may write
\[
  \Lc X = \ul{\Spec}_{\Osc_X} (\Osc_X \otimes_{\Osc_X \otimes \Osc_X} \Osc_X)
\]
where all operations here are implicitly derived.

\subsection{Hochschild homology}

Suppose $X = \Spec R$ is affine.
Then 
\[
  \Osc(\Lc X) = R \otimes_{R \otimes R} R = HC_*(R)
\]
the \emph{Hochschild chains} of $R$.
The \emph{Hochschild homology} $HH^*(R)$ is the homology of this complex.
In general, Hochschild chains are a ``derived version of the abelianization $R / [R, R]$''

\begin{rmk}
  When $X$ is a topological space / Betti stack, we have $\Osc(X) = C^*(X)$.
  The above discussion roughly translates into something like $H^*(\Lc X) = HH_*(C^*(X))$, at least when $X$ is simply connected (so the diagonal is affine).
\end{rmk}

\begin{thm}[Hochschild--Kostant--Rosenberg]
  Let $R$ be a commutative algebra with $X = \Spec R$ smooth.
  Then there is an equivalence of dg-algebras $HC_*(R) \simeq \Omega^*_R = \Sym^* \Omega^1_R[1]$, the complex of differentials on $R$.
  Geometrically, this translates into $\Osc(\Lc X) \simeq \Omega^*_X$.
\end{thm}

\begin{dfn}
  Let $\TT_X[-1]$ be the odd tangent bundle of $X$, i.e.\ $\ul{\Spec}_{\Osc_X} \Sym^* \Omega^1_X[1]$.
\end{dfn}

HKR can be interpreted as an equivalence $\Lsc X \simeq \TT_X[-1]$ for any smooth scheme $X$.
David pointed out that this equivalence should be thought of as an exponential map $\TT_X[-1] \to \Lsc X$.
For stacks, the above statement of HKR fails.

\begin{ex}
  For $X = BG$, the loop space $\Lc G$ is the adjoint quotient $G / G$.
  However, $\TT_X[-1]$ is $\gfr / G$, which is quite different.
\end{ex}

\section{(2/5) Samuel Mayonnaise -- Continued}

Let $X$ be a smooth scheme.
Recall from last time that $\Lc X = \Hom(S^1, X)$ is the \emph{derived loop space} of $X$.
(Here $\Hom$ indicates internal $\Hom$.)
We may also understand $\Lc X$ as a fiber product $X \times_{X \times X} X$.
If $X = \Spec R$ is affine, then $\Osc(\Lc X) = R \otimes_{R \otimes R} R = HC_*(R)$, the Hochschild chains of $R$.

\subsection{HKR}

We'll review the statement of HKR again in more detail.

\begin{dfn}
  The \emph{de Rham algebra} of $X$ is $\Omega^*_X = \Sym^*_{\Osc_X} \Omega_X[1]$.
  The \emph{odd tangent bundle} of $X$ is $\TT_X[-1] = \Spec_{\Osc_X} \Omega^*_X$.
\end{dfn}

The object $\TT_X[-1]$ is a derived scheme.
It consists of $X$ together with ``odd tangent vectors.''

\begin{thm}[Hochschild--Kostant--Rosenberg]
  For a smooth scheme $X$, there is an equivalence $\Lc X \simeq \TT_X[-1]$ (as derived schemes over $X$).
  Via this equivalence, the constant loops correspond to the zero section.
\end{thm}

Part of the idea here is that the ``loops'' in a scheme are all infinitesimal.
When we go to stacks, there are more interesting global loops, so the theorem fails.

The statement of HKR is \'etale-local, so we can reduce to proving the theorem for $\AA^n$.
This is not too hard to do!

\begin{ex}
  Let $X = \AA^n = \Spec k[x_1, \dots, x_n]$, and write $R = k[x_1, \dots, x_n]$.
  We want to compute $HH_*(R) = R \otimes_{R \otimes R} R$.
  There is a \emph{Koszul resolution}
  \[
    \begin{tikzcd}
      0 \rar & \wedge^n (R \otimes R)^{\oplus n} \rar & \dots \rar & \wedge^1 (R \otimes R)^{\oplus n} \rar & R \otimes R \rar & 0
    \end{tikzcd}
  \]
  of $R$, where:
  \begin{itemize}
    \item the resolution map $R \otimes R \to R$ is by $x_i \otimes 1 \mapsto x_i$ and $1 \otimes x_i \mapsto x_i$,
    \item the map $\wedge^1 (R \otimes R)^{\oplus n} \to R \otimes R$ sends the $i$th standard basis vector to $x_i \otimes 1 - 1 \otimes x_i$, and
    \item other maps are defined using alternating sums generalizing the above.
  \end{itemize}
  Let's call this complex $K^*$ and compute $K^* \otimes_{R \otimes R} R$.
  We get
  \[
    \begin{tikzcd}
      0 \rar & \wedge^n R^{\oplus n} \rar & \dots \rar & \wedge^1 R^{\oplus n} \rar & R \otimes R \rar & R
    \end{tikzcd}
  \]
  where all differentials are $0$.
  Then $HH_i(\AA^n) = \wedge^i R^n = \Omega^i_{\AA^n}$, and this upgrades to an isomorphism of complexes $HC_*(\AA^n) \simeq \Omega^{-*}_{\AA^n}$.
  Thus HKR holds for $X = \AA^n$.

  David and Yuji worked out on the board that, when we work on $\AA^1$, the Koszul complex is
  \[
    \begin{tikzcd}
      \Isc_\Delta \rar & \Osc_{\AA^1 \times \AA^1}.
    \end{tikzcd}
  \]
  Tensoring with $\Osc_{\Delta}$ gives 
  \[
    \begin{tikzcd}
      \Isc_\Delta / \Isc_{\Delta}^2 \rar["0"] & \Osc_{\Delta}.
    \end{tikzcd}
  \]
  Using $\Isc_\Delta / \Isc_{\Delta}^2 = \Omega_X$, we get a ``coordinate--free'' proof of HKR in this case.
\end{ex}

\subsection{Loop rotation}

So far we have not encountered the de Rham differential.
Where does this come from?
How can we see it from the loop space perspective?

There's a natural action of $S^1 \curvearrowright \Lc X$ coming from the rotation action $S^1 \curvearrowright S^1$.
This makes sense as an action of ``group objects'' or at the level of functors of points.
Because $S^1$ is connected, every point of $S^1$ acts by ``the identity'' on $\Lc X$, but the action has interesting monodromy (by moving around $S^1$).
In other words, the action corresponds to an interesting self-homotopy of the identity map $\id: \Lc X \to \Lc X$.
To linearize this action, let $\Lambda = C_*(S^1) = k[\epsilon] / \epsilon^2$ (where $\deg \epsilon = -1$).
Then $S^1 \curvearrowright \Lc X$ corresponds to a $\Lambda$-module structure on $\Osc(\Lc X)$.

Let's think about what this means.
Here $\epsilon$ acts on $\Osc(\Lc X)$ by degree $-1$.
We have $\epsilon^2 = 0$ (since $S^1$ has no higher homology).
The fact that $S^1 \curvearrowright \Lc X$ is a group action corresponds to the Leibniz rule $\epsilon (fg) = (\epsilon f) g + f (\epsilon g)$.
In other words, $\epsilon$ serves as a differential making $\Osc(\Lc X)$ into a cdga!

\begin{thm}
  Under the HKR isomorphism, the loop rotation action $S^1 \curvearrowright \Lc X$ corresponds to the de Rham differential $d$ on $\Omega^*_X$.
\end{thm}

\section{(1/12) Will Fisher -- Tate Objects and $D$-Modules}

Our goal for today is to discuss an equivalence
\[
  \IndCoh(\Lc X)^\Tate \simeq \Dsc_X\dMod.
\]
Here $X$ is a (not necessarily smooth) scheme.

\begin{rmk}
  We may understand the categories appearing here as the periodic cyclic homology of $2\IndCoh(X)$.
\end{rmk}

\begin{rmk}
  Recall that $\Lc X = X \times_{X \times X} X$.
  If $X$ is smooth, then $\Lc X$ is ``quasi-smooth,'' and the homological algebra here behaves nicely.
  ``Quasi-smooth'' may be understood as a derived analogue of the lci condition.

  If $\Lc X$ is smooth with smooth diagonal, then $\Lc X$ is smooth.
  This is typically a stacky phenomenon, i.e.\ $\Lc BG = G / G$ is smooth.
\end{rmk}

\subsection{$S^1$-invariants in HKR}

In the previous talks we learned that, when $X$ is a smooth scheme, there is an equivalence $\Lc X \simeq \TT_X[-1] = \Spec_{\Osc_X} \Omega_X^{-*}$.
There is a loop rotation action of $S^1$ on $\Lc X$.
Under the above equivalence, this corresponds to the action of $C_{-*}(S^1) = k[\lambda]$ on $\Omega_X^{-*}$ induced by the de Rham differential.
(Here $|\lambda = -1$, so $\lambda^2 = 0$.)
How can we upgrade the HKR equivalence $\QCoh(\Lc X) \simeq \Omega_X^{-*}\dMod$ to incorporate this $C_{-*}(S^1)$-action?

Let $(C^*, d)$ be a dg $k$-module with $S^1$-action, i.e.\ $(C^*, d) \in \QCoh(BS^1)$.
Descent along $\pt \to BS^1$ (i.e.\ comonadic Barr-Beck) gives an equivalence $\QCoh(BS^1) \simeq \Osc(S^1)\dComod \simeq C^*(S^1)\dComod$.
By $k$-linear duality (or descent using $!$-functors), we may identify this further with $C_{-*}(S^1)\dMod \simeq k[\lambda]\dMod$.
That is, we have an equivalence
\[
  \QCoh(BS^1) \simeq k[\lambda]\dMod.
\]
Taking invariants of objects of $\QCoh(BS^1)$ produces modules over $\Osc(BS^1) = C^*(BS^1) = k[u]$ (with $|u| = 2$).
For $k[\lambda] / \lambda^2$-modules, ``taking invariants'' corresponds to taking $\Hom$ out of the trivial representation $k$.
The resulting complexes have an action of $\Hom_{k[\lambda]}(k, k)$.
It is a good exercise to compute that $\Hom_{k[\lambda]}(k, k) = C^*(BS^1) = k[u]$.

\subsection{Why work with $\IndCoh$?}

When we pass from $k[\lambda]\dMod$ to $k[u]\dMod$, we send $k[\lambda]$ to $k$.
Every object of $k[\lambda]\dMod$ can be built as a colimit of copies of $k[\lambda]$.
Thus the images of such objects are always torsion modules over $k[u]$.
We can actually see that $k[\lambda]\dMod$ can be identified with the (colimit-completion of the) category of $u$-torsion $k[u]$-modules.

We can fix things by working with ind-coherent modules over $k[\lambda]$.
The category $\IndCoh(k[\lambda])$ is compactly generated by $k$.
Thus we get $\IndCoh(k[\lambda]) \simeq \End_{k[\lambda]}(k)\dMod \simeq k[u]\dMod$.
The Tate construction will serve as a generalization of this.

\begin{ex}
  The $k[u]$-module $k[u, u\inv]$ corresponds to the ind-coherent $k[\lambda]$-module built out of infinitely many shifts of $k[\lambda]$.
  Here the shifts are cooked up so that we have one copy of $k$ (or $k \cdot \lambda$) in each degree, with the $k[\lambda]$-module structure sending odd degrees to the next even degrees.
  The ``underlying module'' of this is trivial, but it is still interesting as an ind-coherent module!
\end{ex}

\section{(2/19) Will Fisher -- Continued}

The speaker is posting his own notes online, so look at those for an alternative viewpoint.

Our goal is to show the following: for $X$ a smooth underived scheme, we have
\[
  \IndCoh(\Lc X)^\Tate \simeq \Dc_X[u, u\inv]\dMod.
\]
Here $u$ is a formal variable of degree $2$.

\subsection{The Tate construction}

Let $\Cc$ be a $k$-linear category with $S^1$-action.
Then $\Cc^{S^1}$ is a module over $\Mod_k^{S^1}$ (considered with respect to the trivial $S^1$-action).
We can write $\Mod_k^{S^1} = \QC(\pt)^{S^1} = \QC(BS^1)$.
This last category is $k[[u]]$-linear where $|u| = 2$, so $\Cc^{S^1}$ must also be $k[[u]]$-linear.
We would like to ``invert $u$ on each $\Hom$-space.''

\begin{dfn}
  The \emph{Tate construction} applied to $\Cc$ is
  \[
    \Cc^\Tate = \Cc^{S^1} \otimes_{k[[u]]} k((u)).
  \]
\end{dfn}

\begin{rmk}
  The $0$-categorical analogue of the Tate construction is as follows.
  Take a chain complex $C^\bullet$ with $S^1$-action.
  Note that $C_{-\bullet}(S^1) = k[\lambda]$, so $(C^\bullet)^{S^1} = \Hom_{k[\lambda]}(k, C^\bullet)$.
  The result is linear over $k[[u]] = \End_{k[\lambda]}(k)$, and we can invert $u$ to get the Tate construction $C_{-\bullet}^\Tate$.

  If $X$ is a reasonable compact space with $S^1$-action, then $C^\bullet(X)^{S^1} = C^\bullet_{S^1}(X)$.
  We can show that $C^\bullet(X)^\Tate = C^\bullet(X^{S^1})[u, u\inv]$.
\end{rmk}

\begin{rmk}
  Given a variety $X$ with a function to $\GG_m$, we get an $S^1$-action on $\QC(X)$.
  The Tate construction here gives a definition of the corresponding matrix factorization category.
\end{rmk}

Here's another way to define a ``Tate construction.''

\begin{dfn}
  Let $\Cc$ be a compactly generated $k$-linear category, i.e.\ $\Cc = \Indsf(\Cc^\omega)$.
  Suppose there is an $S^1$-action on $\Cc^{\omega}$ inducing $S^1 \curvearrowright \Cc$.
  Then define $\Cc^{\omega S^1} = \Indsf((\Cc^\omega)^{S^1})$, and let the \emph{$\omega$-Tate construction} be
  \[
    \Cc^{\omega \Tate} = \Cc^{\omega S^1} \otimes_{k[[u]]} k((u)).
  \]
\end{dfn}

Note that $\Indsf((\Cc^\omega)^{S^1}) \neq \Indsf(\Cc^\omega)^{S^1}$ in general, so the $\omega$-Tate and Tate constructions do not agree in general.
We'll see this in the next section.

\subsection{The case $X = \pt$}

To work out what happens for $X = \pt$, we need to understand the $k[[u]]$-linearity on $\QC(\pt)^{S^1} = \QC(BS^1)$.
The affinization map $\pi: BS^1 \to \Aff(BS^1)$ gives an adjunction
\[
  \pi^*: \QC(\Aff(BS^1)) \rightleftarrows \QC(BS^1) : \pi_*.
\]
We may rewrite this adjunction as
\[
  \Osc(BS^1)\dMod \rightleftarrows k[\lambda]\dMod : \Hom_{k[\lambda]}(k, -).
\]
The object $k \in k[\lambda]\dMod$ is not compact, so $\pi_*$ doesn't preserve colimits.
In particular, this adjunction is not an equivalence.

Koszul duality tells us that $\pi^*$ and $\pi_*$ do restrict to equivalences $\Coh(BS^1) \simeq \Perf(k[[u]])$.
Equivalently, we get $\Indsf(\Coh(BS^1)) \simeq \Indsf\Perf(k[[u]]) = k[[u]]\dMod$.\footnote{There is an unfortunate subtlety here:
$\IndCoh(BS^1) \neq \Indsf(\Coh(BS^1))$.
Indeed, the Gaitsgory-type definition of $\IndCoh(BS^1)$ comes from descent, not ``taking $\Indsf$ of $\Coh$.''
In fact, $\IndCoh(BS^1) = \QCoh(BS^1)$ by our above descent argument.
We will ignore these issues and instead focus on $\Indsf(\Coh(BS^1))$, which doesn't satisfy descent.}

\begin{rmk}
  Here's another way of understand what's going on.
  The category $k[[u]]\dMod$ has two ``basic'' compact objects: $k[[u]]$ and $k$.
  The former object corresponds to the compact object $k[\lambda] \in k[\lambda]\dMod$.
  The latter corresponds to $k \in k[\lambda]\dMod$, which is not compact!
  However, if we replace $k[\lambda]\dMod$ by $\Indsf(\Coh(BS^1))$, we are able to match up our compact objects.

  This is analogous to the difference between $S^1 \curvearrowright \pt$ and $S^1 \curvearrowright S^\infty$ in genuine equivariant homotopy theory.
  We can understand this difference by looking at the $S^1$-fixed points of each.
  The differences between $k[\lambda] \curvearrowright C^\bullet(\pt)$ and $k[\lambda] \curvearrowright C^\bullet(S^\infty)$ becomes more transparent in the world of $\Indsf$-coherent sheaves.
\end{rmk}

Let's think about the equivalence $\Indsf(\Coh(BS^1)) \simeq k[[u]]\dMod$.
The induced map $\QCoh(BS^1) \to k[[u]]\dMod$ is \emph{not} given by the usual pushforward.
Instead, we have to resolve a quasicoherent sheaf by perfect complexes and then push forward each perfect complex.
The category $\Perf(BS^1) = \Perf(k[\lambda])$ is generated by $k[\lambda]$, and $\pi_* k[\lambda] = \Hom_{k[\lambda]}(k, k[\lambda]) = k[[u]] / u$.
Thus our equivalence sends $\QCoh(BS^1)$ to the category $(k[[u]]\dMod)^{u\textrm{-nil}}$ of filtered colimits of torsion modules.
This gives our desired description of the $k[[u]]$-linear structure.

\begin{cor}
  $\IndCoh(\pt)^\Tate = \QCoh(\pt)^\Tate = 0$.
\end{cor}

When we work with the $\omega$-Tate construction, we get a more interesting answer.

\begin{cor}
  $\IndCoh(\pt)^{\omega\Tate} = k((u))\dMod$.
\end{cor}

\begin{proof}
  We have 
  \begin{align*}
    \IndCoh(\pt)^{\omega S^1} &= \Indsf(\Coh(\pt)^{S^1}) \\
                              &= \Indsf(\Vect_{k,\textrm{fg}}^{S^1}) \\
                              &= \Indsf(\Coh(BS^1)) \\
                              &\simeq k[[u]]\dMod.
  \end{align*}
  Tensoring over $k[[u]]$ with $k((u))$ gives $k((u))\dMod$.
\end{proof}

\begin{rmk}
  The big picture here: we started with ``na\"ive $S^1$-equivariance'' and then passed to ``genuine $S^1$-equivariance'' (by allowing an ``$S^1$-fixed theory'').
  The Tate construction kills the original na\"ive $S^1$-equivariant theory and only remembers the ``extra information'' governing the fixed points.
\end{rmk}

\section{(2/26) Will Fisher -- Continued}

Recall that our goal is to show $\Indsf(\Coh(\Lc X))^\Tate = \Dsc_X\dMod \otimes_k k[u, u\inv]$.
We previously showed a $k[[u]]$-linear equivalence $\Indsf(\pi_*|_\Coh) = \Indsf(\Hom_{k[\lambda]}(k, -))$ from $\Indsf(\Coh(k[\lambda]\dMod))$ to $k[[u]]\dMod$.

\subsection{The local case}

Let $R$ be a $k$-algebra with $S^1$-action, i.e.\ an algebra object in $(k[\lambda]\dMod, \star)$.
Suppose that $R$ is coherent as a $k$-module.\footnote{Otherwise one needs to use $t$-structures to fix things. See work of Preygel.}
Then $S^1 \curvearrowright R$ gives an action $S^1 \curvearrowright \Coh(R)$.
We may write $\Coh(R) = \Mod_R(\Coh(k))$ and
\[
  \Coh(R)^{S^1} = \Coh(R \rtimes k[\lambda]) = \Mod_R(\Coh(k[\lambda]))
\]
where $R \rtimes k[\lambda] = R\angles{\lambda} / (\lambda^2, [\lambda, r] = \lambda \cdot r)$.
Koszul duality gives a $k[[u]]$-linear equivalence $\Coh(R)^{S^1} \simto R^{S^1}\dMod(\Perf(k[[u]]))$.
Now we need to understand how to compute $R^{S^1} = \Hom_{k[\lambda]}(k, R)$.

Suppose $R = (C^\bullet, d)$.
Using the resolution of $k$ by shifts of $k[\lambda]$ presented in Swapnil's talks, we get $R^{S^1} = (C^\bullet[[u]], d + u\lambda)$.
Thus $\Coh(R)^{S^1} = \Coh(R^{S^1}) = \Coh((C^\bullet[[u]], d + u\lambda))$.

\begin{rmk}
  David stressed that the expression $u\lambda$ involves $u \in k[[u]] = C^\bullet(BS^1; k)$ and $\lambda \in k[\lambda] = C_\bullet(S^1; k)$.
  These live in Koszul dual spaces, and combining them gives something ``dimensionless'' / ``natural'' (in the same way that $d$ is).
\end{rmk}

\subsection{Proving the equivalence}

Let's assume that our story works over a more general base.
Let $X$ be a scheme, and consider the map $p: \Lc X \to X$.
Note that $p$ is affine and $\Osc_{\Lc X}$ is coherent as an $\Osc_X$-module because
\[
  \Lsc X = \TT_X[-1] = \Spec_{\Osc_X} \Omega_X^{-\bullet},
\]
Working locally and using the above results we get
\[
  \Coh(\Lsc X)^{S^1} = (\Omega_X^{-\bullet})^{S^1}\dMod(\Coh(X)) = (\Omega_X^{-\bullet}[[u]], u d_\dR)\dMod(\Coh(X)).
\]
When we invert $u$, we get
\[
  \Coh(\Lsc X)^{S^1} \otimes_{k[[u]]} k((u)) = (\Omega_X^{-\bullet}, d_\dR)\dMod(\Coh(X)) \otimes_k k((u)).
\]
We'll call objects of $(\Omega_X^{-\bullet}, d_\dR)\dMod(\Coh(X))$ (coherent) $\Omega$-modules.
We can relate these to $\Dsc$-modules using $\Omega$-$\Dsc$ duality.
The moral is that, for a Lie algebra $\gfr$, Koszul duality relates modules over the Chevalley--Eilenberg algebra $\wedge^\bullet \gfr^*$ to modules over the universal enveloping algebra $\Uc(\gfr)$.

\begin{rmk}
  David gave a general idea for Koszul duality (originally due to Quillen): the classifying space of a based loop group is a formal neighborhood of the basepoint.
  In symbols, for $x \in X$, we should have $\hat{X}_x = B \Omega_x X$.
  If we specify $y \in \hat{X}_x$ which is ``infinitesimally close'' to $x$, then the $\Omega_x$-torsor of paths from $x$ to $y$ is the universal $\Omega_x$-torsor.

  The above discussion fits into this paradigm: we can think of an odd shift of $\gfr$ as $\hat{X}_x$.
  Taking $\End$ of the augmentation module gives $\Sym \gfr$.
  Turning on the Chevalley--Eilenberg differential (on $\wedge^\bullet \gfr^*$) deforms $\Sym \gfr$ to $\Uc \gfr$.
\end{rmk}

Following the above moral, we can think of $(\Omega^{-\bullet}_X, d_\dR) = \wedge^{-\lambda} \Tsc_X^\vee$, so we get a Koszul duality with $\Dsc_X = \Uc(\Tsc_X)$.
We note that $\scEnd_{(\Omega_X^{-\bullet}, d_\dR)}(\Osc_X) = \Dsc_X$.
Koszul duality gives 
\[
  (\Omega_X^{-\bullet}, d_\dR)\dMod(\Coh(X)) = \Perf(\Dsc_X^\bullet).
\]
Taking $\Indsf$ of both sides gives
\[
  \Indsf(\Coh(\Lsc X)^{S^1}) \otimes_{k[[u]]} k((u)) = \Dsc_X\dMod \otimes_k k((u)).
\]

\begin{rmk}
  David gave the following teaser for next time.
  So far, given a scheme $X$, we have looked at $\TT_X[-1]$ with the de Rham differential.
  This is Koszul dual to $T^*_X[2]$ with its symplectic form $\omega$.

  More generally, for a stack $X$, we should think of $T_X[-1]$ as the Lie algebra of the loop space $\Lsc X$.
  We don't know what the correct dual to this is!
  Potent categorical representations should help us understand things away from the ``identity section'' of $\Lsc X$.
\end{rmk}

\section{(3/5) Jacob Erlikhman -- Potent $D$-Modules versus Ordinary $D$-Modules}

Our goal is to discuss how some of the above results break down for stacks.
We shall begin by reviewing some basic stack theory.

\subsection{Classifying stacks}

The classifying stack $BG$ can be constructed as the geometric realization of the simplicial diagram 
\[
  \begin{tikzcd}
    \dots G^2 \rar[shift left] \rar \rar[shift right] & G \rar[shift left] \rar[shift right] & \pt.
  \end{tikzcd}
\]
The maps here come from projections and multiplication maps.
We can also think of $BG$ as the sheafification of the na\"ive quotient $S \mapsto \pt(S) / G(S)$.
Given a map $\gamma: S \to BG$, pulling back $\pt \to BG$ gives a principal $G$-bundle $P \to S$.
Thus maps to $BG$ classify (and are classified by) principal $G$-bundles.

\subsection{Quotient stacks}

Quotient stacks $X / G$ behave similarly.
For instance, they are the geometric realization of the simplicial diagram 
\[
  \begin{tikzcd}
    \dots X \times G^2 \rar[shift left] \rar \rar[shift right] & X \times G \rar[shift left] \rar[shift right] & X.
  \end{tikzcd}
\]
We may also view $X / G$ as the sheafification of $S \mapsto X(S) / G(S)$.

\begin{rmk}
  One can construct a simplicial diagram $X^\bullet$ for any space $X$ (using the projection maps as natural maps).
  If we replace some factors by $G$ and quotient the whole thing out by $G$, we get the diagram mentioned above.
  David suggested thinking of this as a ``free $G$-resolution of $X$.''
\end{rmk}

\begin{ex}
  If $P$ is a principal $G$-bundle on $X$, then $P = S \times_{BG} \pt$.
  We can compute
  \[
    P / G = S \times_{BG} \colim_{\Delta\op} G^n = S \times_{BG} BG = S.
  \]
  That is, quotienting a $G$-torsor by $G$ gives the underlying space.
\end{ex}

\begin{ex}
  For $H \subset G$ a subgroup, we have $G \bs (G / H) = BH$.
\end{ex}

Given any $S \to X / G$, we get an induced map $S \to BG$, and hence a principal $G$-bundle $P \to S$.
We let $P \times^G X = S \times_{BG} X / G$.
We can construct a cube
\[ % taken from TikZcd documentation
\begin{tikzcd}
& X \times^G P \arrow[rr] \arrow[dd] & & X / G \arrow[dd] \\
  X \times P \ar[ur] \arrow[rr, crossing over] \arrow[dd] & & X \ar[ur] \\
& S \arrow[rr] & & BG \\
  P \arrow[rr] \ar[ur] & & \pt \arrow[from=uu, crossing over] \ar[ur] \\
\end{tikzcd}
\]
and show that $X \times^G P = (X \times P) / G$.
By studying this cube we see that $G$-torsors on $S$ correspond to $G$-torsors $P$ on $S$ together with:
\begin{itemize}
  \item a section of $X \times^G P \to S$, or equivalently:
  \item a $G$-equivariant map $P \to X$.
\end{itemize}

\begin{ex}
  Let $X = \AA^1$ and $G = \GG_m$.
  Then $(\AA^1 / \GG_m)(S)$ parametrizes line bundles $\Lsc$ on $S$ together with a choice of section $\sigma \in H^0(X, \Lsc)$.
  The data of the section may also be understood as a $\GG_m$-equivariant map from $\Lsc^\times$ to $\AA^1$.
  Roughly, $\sigma$ corresponds to $s \mapsto [\ell_s, \phi(\ell_s)] \in (\Lsc^\times) \times^{\GG_m} \AA^1$, and the $\GG_m$-equivariant map is $\phi$.
\end{ex}

\subsection{Loops on $BG$}

Let's compute $\Lc BG = \Map(S^1, BG)$.
This has $S$-points $\Map(S \times S^1, BG)$.
Because $S^1 = B\ZZ$, we have $S \times S^1 = S / \ZZ$ (where the $\ZZ$-action on $S$ is trivial).
Thus $S$-points of $\Lc BG$ correspond to $G$-torsors on $S / \ZZ$, or equivalently $\ZZ$-equivariant $G$-torsors $P$ on $S$.
The data of each such object corresponds to a $G$-torsor on $S$ together with a choice of automorphism $a: P \to P$.

Now let's consider the adjoint quotient $G / G$.
This has $S$-points given by $G$-torsors $P$ on $S$ together with a $G$-equivariant map $f: P \to G$.
Given $a \in \Aut(P)$, we may define $f: P \to G$ via $a(p) = p f(p)$.
One checks that $a$ is a bundle automorphism if and only if $f$ is $G$-equivariant for the adjoint action.
This construction is invertible, so we get an identification $\Lc BG = G / G$.

\section{(3/12) Jacob Erlikhman -- Continued}

\subsection{Review of $\Lc BG$}

Last time we talked about $\Lc BG \simeq G / G$.
Recall:
\begin{itemize}
  \item $\Lc BG$ has $S$-points given by $G$-torsors $P$ on $S$ together with $a \in \Aut(P)$.
    Here $S^1$ acts on $(\Lc BG)(S)$ by $a$ (thought of as an automorphism of $\id_{(P,a)}$).
  \item $G / G$ has $S$-points given by $G$-torsors $P$ on $S$ together with $G$-equivariant maps $P \to G$.
    The $S^1$-action here is by the distinguished automorphism $a$, defined by $a(p) = p \cdot f(p)$.
\end{itemize}
When $G = T$ is commutative (e.g.\ a torus), we have $\Lc BT = T \times BT$.
\begin{itemize}
  \item $T \times BT$ has $S$-points given by $T$-torsors $P$ on $S$ together with $S$-points of $T$.
    The $S^1$-action is again by the distinguished automorphism $a$.
\end{itemize}

\begin{rmk}
  Peisheng asked about the relationship between $G / G$ and $T \sslash W$.
  David explained that we have $\Osc(G / G) = \Osc(G)^G = \Osc(T)^W = \Osc(T \sslash W)$.
  There are many relationships between these: for example, there is a characteristic polynomial map $\chi: \gfr / G \to \cfr = \tfr \sslash W$.
  This has a \emph{Kostant section} $\kappa: \tfr \sslash W \to \gfr / G$.
\end{rmk}

\subsection{The $S^1$-action on $\Coh(G / G)$}

Let's consider the action map $\act: S^1 \times G / G \to G / G$.
Recall that coherent sheaves on $S^1 \times X$ are coherent sheaves on $X$ together with an automorphism.
Pulling back along $\act^*$ therefore sends $\Fsc \in \Coh(G / G)$ to a pair $(\Fsc, a_\Fsc)$ where $a_\Fsc$ is a distinguished automorphism.
What is this?

Suppose we are given an $S$-point $g: S \to G / G$.
Then we get $a_g: \Fsc_g \to \Fsc_g$, defined by the action of $g$.

\begin{ex}
  Let $S = \Spec k$, and let $g: \Spec k \to G / G$.
  Then $\Fsc_g \in \Coh(\Spec k)$ is just a (dg-)vector space.
  The element $g$ centralizes itself, and the action $a_g$ is the induced conjugation action on the stalk.
\end{ex}

\subsection{$S^1$-invariants in $\Coh(T \times BT)$}

\begin{ex}
  We'd like to classify $S^1$-invariant sheaves on $\Osc_{T \times BT}$ for $T = \GG_m$.
  These are sheaves on which, at every $t \in T$, $t$ acts by $1$.
  We can describe these as follows:
  \begin{itemize}
    \item Any pullback of a coherent sheaf on $T$ is $S^1$-invariant.
    \item Any twist $\Osc_1 \otimes \Osc(k)$ is $S^1$-invariant.
    \item Any twist $\Osc_{\zeta_k} \otimes \Osc(k)$ (with $\zeta_k$ a $k$th root of unity) is $S^1$-invariant.
  \end{itemize}
\end{ex}

A similar (but harder) computation may be performed for $G / G$ with $G$ nonabelian (?).

\begin{ex}
  Let $T$ be an algebraic torus.
  Then we can write
  \[
    \Coh(T \times BT) = \Coh(T) \boxtimes \Rep^\coh(T) = \oplus_{\lambda \in \weight(T)} \Coh(T)_\lambda.
  \]
  We can compute $S^1$-invariants and get
  \begin{itemize}
    \item $\Coh(T)_\lambda^{S^1} = \Osc(T) / (t^\lambda - 1)\dMod$ for $\lambda \neq 0$.
      This is secretly $k[\epsilon]$-linear (where $|\epsilon| = -1$) and $\epsilon$ acts essentially by $t^\lambda - 1$ (?).
    \item $\Coh(T)_0^{S^1} = \Osc(T)[\epsilon]\dMod$.
  \end{itemize}
  Both essentially come from computing $\Osc(T) \otimes_{\Osc(\GG_m)} k$.
\end{ex}

Geometrically, what is going on here is that we have a map $X \to \GG_m$ (here $X = T / T$) classifying the action of automorphisms.
Coherent sheaves on the (derived) fiber $X_1 = X \times_{\GG_m} 1$ correspond to coherent sheaves on $X$ with the automorphism identified with $1$.
That is, $\Coh(X_1) = \Coh(X)^{S^1}$.
Now observe that $X_1$ has an action of $1 \times_{\GG_m} 1 = \Spec k[\epsilon]$ (with $|\epsilon| = -1$).
This induces $\Coh(1 \times_{\GG_m} 1) \curvearrowright \Coh(X_1)$, where $\Coh(1 \times_{\GG_m} 1)$ is equipped with the convolution product.
The monoidal unit of $\Coh(1 \times_{\GG_m} 1)$ is the augmentation module $\onebb = \Osc_1$.
We have $\End(\onebb) = k[u]$, so Koszul duality should turn $\Coh(X_1)$ into a $k[u]$-linear category.

\begin{ex}
  We can show that $\Osc(T)[\epsilon]\dMod^\coh$ is Koszul dual to $\Osc(T)[u]\dMod^\coh$.
  Inverting $u$ takes this to $\Osc(T)[u, u\inv]\dMod^\coh$.
  This is $\Coh(T / T)^\Tate$.
  Completing this around $1 \in T$ would give coherent modules over $\Osc(\hat{T})[u, u\inv] = H^*(BT)[u, u\inv]$, which agrees with the category of $\Dsc$-modules on $BT$ (with $u, u\inv$ adjoined).
  Potent categorical representations should somehow capture global data (which is much more subtle than this).
\end{ex}

\section{(4/2) David Nadler -- Potent Categorical Representations}

Let $G$ be a reductive algebraic group over $\CC$ associated with a compact Lie group $G_c$.
We may also define a Langlands dual group $G^\vee$ (over $\CC$).

In his ICM address, Teleman proposed an equivalence between:
\begin{itemize}
  \item (A-model) ``$G_c$-categories.''
    Na\"ively this means objects of $\Loc(G_c)\dMod$ (where the monoidal structure on $\Loc(G_c)$ is convolution).
    We'll be deliberately vague about what conditions (presentability, stability, etc.) we impose on these $\Loc(G_c)$-modules.
  \item (B-model) Rozansky--Witten theory of the holomorphic symplectic group scheme $J_{G^\vee}$.
\end{itemize}
The goal of David's work here (joint with Ben-Zvi and Stefanich) is to give a proper definition of the left hand side.
We'll discuss many possible definitions and the relations between the two.

\subsection{The de Rham story}

Teleman's story is a ``Betti'' story.
Potent categorical representations will a ``de Rham analogue.''

For any algebraic group $G$, we can construct a derived subgroup $G^\der = [G, G] \subset G$ and an abelianization map $a: G \to G^\ab = G / G^\der$.
The left action of $G^\der$ on $G$ gives a moment map $T^* G \to (\gfr^\der)^*$.
Working with $D$-modules, we define $D(G_\ssrm)$ as the full subcategory of $D(G)$ consisting of $D$-modules with singular support above $0 \in (\gfr^\der)^*$.
We think of ``$G_\ssrm$-categories'' as $D(G_\ssrm)$-module categories.

\begin{rmk}
  The group $G_c$ arises if we think of $G$ as a homotopy type.
  Conjecturally, $G_\ssrm$ is what we get if we think of $G$ as an $\AA^1$-homotopy type.
\end{rmk}

\begin{ex}
  Let $T \cong \GG_m^n$ be an algebraic torus, and let $\Lambda = \Hom(\GG_m, T)$.
  Then $\Loc(T_c) = \CC[\Lambda]\dMod = \QC(T^\vee)$.
  Meanwhile, we have $D(T_\ssrm) = D(T) = \QC(\tfr^\vee / \Lambda^\vee)$.
  The two theories here are not the same algebraically, but they do look the same from an analytic perspective.
  Both of these equivalences can be made symmetric monoidal (working with convolution on the left and tensor product on the right).

  Let's understand why $D(T) = \QC(\tfr^\vee / \Lambda^\vee)$, at least for $T = \GG_m = \Spec \CC[z, z\inv]$.
  Here $D$-modules are modules over $D_T = \CC\angles{z, z\inv, z\partial_z}$.
  We think of $\tfr^\vee = \AA^1 = \Spec \CC[z \partial_z]$.
  The operators $z, z\inv$ give $\ZZ$-equivariance on this.
  Under this equivalence, $D_T$ corresponds to $\pi_* \Osc_{\tfr^\vee}$, where $\pi: \tfr^\vee \to \tfr^\vee / \Lambda^\vee$ is the ``universal cover.''

  On the Betti side, we get an equivalence 
  \[
    \Loc(T_c)\dMod = \QC(T^\vee)\dMod = 2\QC(T^\vee).
  \]
  This sits within $\RW(J_{T^\vee})$, where $J_{T^\vee} = T^* T^\vee$.
  Likewise, on the de Rham side, we get 
  \[
    D(T_\ssrm)\dMod = \QC(\tfr^\vee / \Lambda^\vee)\dMod = 2\QC(\tfr^\vee / \Lambda^\vee).
  \]
  This sits within $\RW(J_{T^\vee}^\nabla$, where $J_{T^\vee}^\nabla = T^*(\tfr^\vee / \Lambda^\vee))$.
\end{ex}

In general, we will get an equivalence between:
\begin{itemize}
  \item (A-model) ``$G_\ssrm$-categories.''
    Na\"ively this is $D(G_\ssrm)\dMod$.
  \item (B-model) Rozansky--Witten theory of $J_{G^\vee}^\nabla$.
    Here $J_{G^\vee}^\nabla$ is the ``de Rham group scheme of regular centralizers.''
\end{itemize}

\begin{rmk}
  To see more of these stories, we should enhance both sides.
  For example, we can replace $2\QC(J_{G^\vee}^\nabla$ by $2\IndCoh(J_{G^\vee}^\nabla)$.
  Maybe this is too large for our purposes, maybe not.
\end{rmk}

\subsection{What is $J_{G^\vee}$?}

David first encountered $J_{G^\vee}$ in Ng\^o's proof of the Fundamental Lemma.
Drinfeld has some handwritten notes on this subject which David says may be the best reference for these.\footnote{The notes can be found at \url{https://math.uchicago.edu/~drinfeld/langlands/Regular_centralizers.pdf} -- thanks to Sam Mayo for finding the reference.}

At the beginning of time we have to make a choice: do we study $\gfr^\vee / G^\vee$ or $(\gfr^\vee)^* / G^\vee$?
Drinfeld focuses on the former, while we focus on the latter.
We can think of this as a shifted cotangent bundle.

Chevalley tells us that the invariant-theoretic quotient here is $(\hfr^\vee)^* \sslash W \cong \hfr \sslash W = \cfr_\gfr$.
This is the \emph{Chevalley base} for $\gfr$, and there's a map $\chi: (\gfr^\vee)^* / G^\vee \to \cfr_\gfr$.
We also have a Kostant section $\cfr_\gfr \to (\gfr^\vee)^{*,\textrm{reg}} / G^\vee \subset (\gfr^\vee)^* / G^\vee$
The stack $(\gfr^\vee)^{*,\textrm{reg}} / G^\vee$ is the classifying stack $BJ_{G^\vee}$.
In other words, $J_{G^\vee}$ is the stabilizer of the Kostant section.

It is a good exercise to compute this all for $\SL_2$.

\section{(4/16) David Nadler -- $J$-Day}

Let $G$ be a reductive group over $\CC$ and $\gfr$ its Lie algebra.
We like to study the adjoint representation, encoded by the quotient stack $\gfr / G = T[-1] BG$.
This is equivalent to $\gfr^* / G = T^*[1] BG$ via a choice of inner product.
However, we will keep the two distinct to keep track of ``tangent versus cotangent'' differences.
We'll focus on the $\gfr^* / G$ perspective, though Drinfeld's notes focus on $\gfr / G$.

\subsection{The definition of $J$}

The affinization of $\gfr^* / G$ is the \emph{Chevalley base} $\cfr_{\gfr^*} = \gfr^* \sslash G$.
A theorem of Chevalley identifies $\gfr^* \sslash G = \hfr^* \sslash W$.
The affinization map $\chi: \gfr^* / G \to \cfr_{\gfr^*}$ can be thought of as a characteristic polynomial map.
There's a \emph{Kostant section} $\kappa: \cfr_{\gfr^*} \to (\gfr^*)^\reg$ (landing in the regular locus).
We define
\[
  J_G = \Stab(\kappa)
\]
as the group scheme of regular centralizers (this is a smooth abelian group scheme over $\cfr_{\gfr^*}$).
We can think of the map $J_G \to \cfr_{\gfr^*}$ as a symplectic integrable system.
We may also write $(\gfr^*)^\reg / G = BJ_G$.

\subsection{Examples}

\begin{ex}
  Let $G = T$ be a torus.
  Then $\tfr^* / T = \tfr^* \times BT$, and the map $\chi: \tfr^* \times BT \to \tfr^*$ is just the projection.
  The Kostant section here is the identity.
  Thus $J = \tfr^* \times T$ is the cotangent bundle of $T$.
\end{ex}

\begin{ex}
  Let $G = \SL_2$.
  The map $\chi: \slfr_2^* / G \to \cfr_{\slfr_2^*}$ is the determinant:
  \[
    \chi \begin{pmatrix}
      a & b \\
      c & -a
    \end{pmatrix} = -a^2 - bc.
  \]
  Every conjugacy class in $\slfr_2^*$ other than $0$ is regular.
  The Kostant section (which depends on a choice of Borel) is given by matrices in rational normal form:
  \[
    \kappa(\lambda) = \begin{pmatrix}
      0 & -1 \\
      \lambda & 0
    \end{pmatrix}.
  \]
  Note that we do not want to use diagonal matrices here (there are two diagonal matrices in $\slfr_2^*$ for any choice of determinant, so we don't get a well-defined section).
  We can think of the image of the Kostant section as a translation of $\bfr^\bot$.

  The stabilizer over a regular semisimple element of the Kostant section has two components, with one given by the maximal torus $T \subset \SL_2$.
  The stabilizer of the nilpotent conjugacy class is the group of matrices
  \[
    \begin{pmatrix}
      \pm 1 & u \\
      0 & \pm 1
    \end{pmatrix},
  \]
  which we can think of as the center of $\SL_2$ times the unipotent subgroup of $B$.
  
  Note that $\pi_1(\cfr_{\gfr^*} \setminus 0) = \ZZ$, so we can try to understand the monodromy of the fibers $T_\lambda$.
  Let's start by computing each $T_\lambda$: from the centralization equation, we see that the stabilizer consists of matrices
  \[
    \begin{pmatrix}
      a & b \\
      -b\lambda & a
    \end{pmatrix}.
  \]
  with $a^2 + b^2 \lambda = 1$.
  This has eigenvalues $a \pm i \sqrt{\lambda} b$.
  As we move around $0$, we swap the eigenvalues $\sqrt{\lambda}$ and $-\sqrt{\lambda}$.
  We can think of $\tfr^* \setminus 0 \to \cfr_{\gfr^*} \setminus 0$ as a double cover (more generally, a $W$-fold cover over the regular semisimple locus) over which the monodromy trivializes.
  We should think of the central fiber (over $0$) as a degeneration of these.
\end{ex}

\begin{ex}
  For $\GL_n$, consider the action $W \curvearrowright \tfr$.
  For every $\lambda \in \cfr_{\glfr_n^*}$, we get a finite subscheme $\tfr_\lambda \subset \tfr$ (the fiber of $\tfr \to \tfr \sslash W$ over $\lambda$).
  We can write $J|_\lambda = \Osc(\tfr_\lambda)^\times$.
  David claims that this is Cayley--Hamilton in disguise.
  There's a paper of Donagi and Gaitsgory doing a (complicated) version of this for more general groups.

  For $\GL_2$, this looks like the projection of $\mu^2 = \lambda$ onto the $\lambda$ coordinate.
  The fiber over any $\lambda \neq 0$ is two points, so $J_\lambda = \Osc(\tfr_\lambda)^\times = \GG_m \times \GG_m$.
  The fiber over $0$ is $\Spec \CC[\epsilon] / \epsilon^2$, so $J_0$ consists of $a_0 + a_1 \epsilon$ with $a_0 \neq 0$.
\end{ex}

\subsection{The Ng\^{o} homomorphism}

Consider the fiber product $\tilde{J}_G = J_G \times_{\cfr_{\gfr^*}} \gfr^*$.
Let $I_G$ be the set of $(x, g) \in \gfr^* \times G$ with $\Ad_g^* x = x$.

\begin{thm}
  There is a homomorphism of group schemes (over $\gfr^*$) $n: \tilde{J}_G \to I_G$ which is an isomorphism over $(\gfr^*)^\reg$.
\end{thm}

\begin{proof}
  The irregular locus in $\gfr^*$ has codimension at least $3$, so we can extend the identity homomorphism uniquely by working at the level of functions and using Hartogs' theorem.
\end{proof}

The Ng\^{o} homomorphism $n$ gives us a way to control $\gfr^* / G$ using the regular part.
This is useful in relating $J$ to the Hitchin fibration.

\section{(4/30) David Nadler -- Continued}

No seminar next week, but there will be a hike.
David recommends you watch the documentary ``Hands on a Hard Body'' for an idea of what we will be in for.
You can watch David's recent talk at SLMath for more on potent categorical representations.

\subsection{Review}

Let $G$ be a reductive group.
There's a natural characteristic polynomial map $\chi: \gfr^* \to \cfr_\gfr$, and there's a Kostant section $\kappa: \cfr_\gfr \to \gfr^\reg$.
We define $J_G^{\textrm{tangent}}$ as the stabilizer group scheme of $\kappa$.
This is a commutative group scheme over $\cfr_\gfr$.

We can tell a similar story in the cotangent picture, giving $J_G \to \cfr_{\gfr^*}$.
This is a completely integrable system.
The fibers are Lagrangian and generically look like the maximal torus of the group (the map is generically $(T \times \tfr) / W \to \tfr / W$).
As we degenerate to the origin, the fibers get closer and closer to unipotent groupss.

For today we will focus on $J_{\check{G}}$ (the universal centralizer for $(\check{\gfr})^*$), which lives over $\cfr_\gfr$.

\begin{ex}
  For $\check{G} = \SL_2$, we are looking at a group scheme over $\AA^1$ with:
  \begin{itemize}
    \item Generic fiber $\GG_m$, and
    \item Special fiber $Z(\SL_2) \cdot \GG_a$.
  \end{itemize}
  This is an open in $\Stab_{\check{G}}((\check{\gfr})^*)$.
\end{ex}

\subsection{Why do we care?}

Here's a quick explanation of why we should care about $J$.
We care about the geometry of $(\check{\gfr})^*$ with its adjoint action.
This is complicated!
To simplify this, we can throw away the irregular locus.
This locus has codimension $3$, so throwing it away doesn't change the functions on our space.
Furthermore, everything that shows up when we study $J$ is commutative, so things are much simpler.

\subsection{A de Rham version of $J$}

We can think of $J_{\check{G}}$ as a ``Betti-type object'' -- it is based on understanding a sort of monodromy.
We can write $J_{\check{G}} = \Hom(BZ, BJ_{\check{G}})$, where $\Hom$ is abelian group stack homomorphisms.

The ``de Rham version'' is 
\[
  J_{\check{G}}^\nabla = \Hom((\GG_m)_\dR, BJ_{\check{G}}).
\]
This is again a symplectic commutative group scheme over $\cfr_\gfr$.
Generically, $J_{\check{G}}^\nabla$ looks like $(\check{t} / \check{\Lambda} \times (\check{t})^*) / W \to \cfr_{\gfr}$.

\subsection{Automorphic realizations of $J_{\check{G}}$}

Note that 
\[
  \cfr_{\check{\gfr}^*} = \cfr_{\gfr} = \gfr \sslash G = \Spec H^*_G(\pt).
\]
We'd like to understand if there's a related topological construction of $J_{\check{G}}$.

Consider the affine Grassmannian $\Gr_G = G((t)) / G[[t]]$.
The quotient 
\[
  \Gsc_{S^2} = G[[t]] \bs \Gr_G = G[[t]] \bs G((t)) / G[[t]]
\]
is equivalent to $\Bun_G$ of the raviolo $\DD \times_{\DD^\times} \DD$ (which is topologically $S^2$).
Here's the reason we've worked with $\check{G}$ throughout:

\begin{thm}[Bezrukavnikov--Finkelberg--Mirkovic]
  The map $J_{\check{G}} \to \cfr_\gfr$ is can be written as $\Spec$ of the natural map
  \[
    H_G^*(\pt) \to H_*(\Gsc_{S^2})
  \]
  induced by $BG \simeq G[[t]] \bs G[[t]] / G[[t]] \hookrightarrow \Gsc_{S^2}$.
\end{thm}

\begin{rmk}
  One can describe the group structure on $J_{\check{G}}$ using this picture, but we won't do this here.
\end{rmk}

The above is a Betti theorem -- how does the de Rham story work?
Let's assume for simplicity that $\check{G}$ is of adjoint type (or more generally has connected center).
Then $J_{\check{G}}^\nabla = T^* \cfr_\gfr / L$, where $L$ is the ``coweight lattice'' $\Hom(\GG_m, J_{\check{G}}$ (varying with fibers).
That is, $J_{\check{G}}^\nabla$ can be understood entirely in terms of $\cfr_\gfr$ and a (varying) discrete group.

\begin{ex}
  For $\PGL_2$, $L \to \AA^1$ has fiber $\ZZ$ away from $0$ and has fiber $0$ over $0$.
  There is a $-1$ monodromy as we move around $0$.
\end{ex}

\end{document}
